\section*{Hoofdstuk \ref{ch:g12:trigo} --- Goniometrische functies}

\begin{solution}{exo:g12:trigo:1}
$\cos\frac{2\pi}{3} = \cos\left(\pi - \frac\pi3\right) = -\cos\frac\pi3
= -\frac12$.

$\sin\frac{5\pi}{6} = \sin\left(\pi - \frac\pi6\right) = \sin\frac\pi6
= \frac12$.

$\cos\frac{7\pi}{4} = \cos\left(2\pi - \frac{\pi}{4}\right)
= \cos\left(-\frac\pi4\right) = \frac{\sqrt2}{2}$.

$\sin\left(-\frac\pi3\right) = -\sin\frac\pi3 = -\frac{\sqrt3}{2}$.
\end{solution}

\begin{solution}{exo:g12:trigo:2}
$f'(x) = 2\sin x\cos x = \sin 2x$ (kettingregel).

$g'(x) = -3\sin(3x) + \sin x + x\cos x$ (ketting- en productregel).

Quotiëntregel:
\[
h'(x) = \frac{\cos x\,(2 + \cos x) - \sin x\,(-\sin x)}{(2+\cos x)^2}
= \frac{2\cos x + \cos^2 x + \sin^2 x}{(2+\cos x)^2}
= \frac{2\cos x + 1}{(2+\cos x)^2}.
\]
\end{solution}

\begin{solution}{exo:g12:trigo:3}
\emph{(a)} Een bijzondere oplossing van $\sin x = \frac{\sqrt3}{2}$ is
$\frac\pi3$. Algemene oplossingen: $x = \frac\pi3 + 2k\pi$ of
$x = \pi - \frac\pi3 + 2k\pi = \frac{2\pi}{3} + 2k\pi$. In
$\intoc{-\pi}{\pi}$: $x \in \left\{\frac\pi3, \frac{2\pi}{3}\right\}$.

\emph{(b)} $\cos 2x = \frac{\sqrt2}{2} = \cos\frac\pi4$ geeft
$2x = \pm\frac\pi4 + 2k\pi$, dus $x = \pm\frac\pi8 + k\pi$. In
$\intoc{-\pi}{\pi}$:
$x \in \left\{-\frac{7\pi}{8}, -\frac\pi8, \frac\pi8,
\frac{7\pi}{8}\right\}$.
\end{solution}

\begin{solution}{exo:g12:trigo:4}
Met $a = \frac\pi3$ en $b = \frac\pi4$:
\[
\cos\frac{\pi}{12} = \cos a\cos b + \sin a \sin b
= \frac12\cdot\frac{\sqrt2}{2} + \frac{\sqrt3}{2}\cdot\frac{\sqrt2}{2}
= \frac{\sqrt2 + \sqrt6}{4},
\]
\[
\sin\frac{\pi}{12} = \sin a\cos b - \cos a \sin b
= \frac{\sqrt3}{2}\cdot\frac{\sqrt2}{2} - \frac12\cdot\frac{\sqrt2}{2}
= \frac{\sqrt6 - \sqrt2}{4}.
\]
\end{solution}

\begin{solution}{exo:g12:trigo:5}
Stel $s = \sin x$: $2s^2 - s - 1 = (2s + 1)(s - 1) \geq 0$ precies wanneer
$s \leq -\frac12$ of $s = 1$.

$\sin x = 1$ geeft $x = \frac\pi2$. $\sin x \leq -\frac12$: op de
eenheidscirkel is dat de boog onder de horizontale rechte $y = -\frac12$, wat
in $\intcc{0}{2\pi}$ overeenkomt met
$x \in \intcc{\frac{7\pi}{6}}{\frac{11\pi}{6}}$. Bijgevolg is
\[
S = \left\{\frac{\pi}{2}\right\} \cup
\intcc{\frac{7\pi}{6}}{\frac{11\pi}{6}} .
\]
\end{solution}

\begin{solution}{exo:g12:trigo:6}
$\dfrac{\sin 3x}{x} = 3\,\dfrac{\sin 3x}{3x} \to 3 \cdot 1 = 3$
(samenstelling met $u = 3x \to 0$).

Met $1 - \cos x = \dfrac{\sin^2 x}{1 + \cos x}$:
\[
\frac{1 - \cos x}{x^2}
= \left(\frac{\sin x}{x}\right)^{\!2} \frac{1}{1 + \cos x}
\xrightarrow[x\to0]{} 1 \cdot \frac12 = \frac12 .
\]

Met de substitutie $u = \frac1x \to 0^+$:
$x \sin\frac1x = \frac{\sin u}{u} \to 1$.
\end{solution}

\begin{solution}{exo:g12:trigo:7}
$f'(x) = 1 - \sin x \geq 0$ voor alle $x$, met gelijkheid precies wanneer
$\sin x = 1$, dus $x = \frac\pi2 + 2k\pi$. Op $\intcc{0}{2\pi}$ stijgt $f$
van $f(0) = 1$ tot $f(2\pi) = 2\pi + 1$ (het nulpunt van $f'$ in
$\frac{\pi}{2}$ is geïsoleerd, zodat de stijging zelfs strikt is); \omterm{def:g10:functions:extrema}{minimum}
$1$ in $0$, \omterm{def:g10:functions:extrema}{maximum} $2\pi + 1$ in $2\pi$. Op $\R$ is $f' \geq 0$ overal met
alleen geïsoleerde nulpunten, dus is $f$ (strikt) \omterm{def:g12:seq:monotonic}{stijgend} volgens
\cref{thm:g12:deriv:variations}, hoewel $f'$ in elke $\frac\pi2 + 2k\pi$
verdwijnt.
\end{solution}

\begin{solution}{exo:g12:trigo:8}
\emph{1.} Zij $R = \sqrt{A^2 + B^2} > 0$. Het punt
$\left(\frac AR, \frac BR\right)$ ligt op de eenheidscirkel, dus bestaat er
een $\varphi$ met $\cos\varphi = \frac AR$ en $\sin\varphi = \frac BR$. Dan
is
\[
R\cos(x - \varphi) = R\bigl(\cos x\cos\varphi + \sin x \sin\varphi\bigr)
= A\cos x + B\sin x = f(x).
\]

\emph{2.} Omdat $\cos$ het \omterm{def:g10:functions:function}{beeld} $\intcc{-1}{1}$ heeft, heeft $f$ \omterm{def:g10:functions:extrema}{maximum} $R$
en \omterm{def:g10:functions:extrema}{minimum} $-R$. Voor $\cos x + \sin x$: $A = B = 1$, $R = \sqrt2$ en
$\varphi = \frac\pi4$, zodat de \omterm{def:g10:algebra:equation}{vergelijking}
$\sqrt2\cos\left(x - \frac\pi4\right) = 1$ wordt, dus
$\cos\left(x - \frac\pi4\right) = \frac{\sqrt2}{2}$, waaruit
$x - \frac\pi4 = \pm\frac\pi4 + 2k\pi$: in $\intoc{-\pi}{\pi}$ is
$x \in \left\{0, \frac\pi2\right\}$.
\end{solution}

\begin{solution}{exo:g12:trigo:9}
$(\sin)'' = (\cos)' = -\sin$ en $(\cos)'' = (-\sin)' = -\cos$: allebei
lossen ze $y'' = -y$ op.

Zij $g = y - a\cos - b\sin$. Dan is $g'' = -g$ (lineariteit van het
afleiden), $g(0) = y(0) - a = 0$ en $g'(0) = y'(0) - b = 0$. Stel
$E = g^2 + (g')^2$. Dan is
\[
E' = 2g g' + 2g' g'' = 2g g' - 2g' g = 0,
\]
dus is $E$ constant en gelijk aan $E(0) = 0$. Een som van twee kwadraten die
identiek verdwijnt, dwingt $g = 0$ af, dus $y = a\cos + b\sin$.
\end{solution}

\begin{solution}{pb:g12:trigo:1}
\textbf{1.} $0.99833\ldots$ en $0.99998\ldots$: kruipend naar $1$.

\textbf{2.} Driehoek $OAM$: basis $1$, hoogte $\sin x$: oppervlakte
$\frac{\sin x}{2}$. Sector $OAM$: het deel $\frac{x}{2\pi}$ van de
eenheidsschijf: oppervlakte $\frac x2$. Driehoek $OAT$: basis $1$, hoogte
$\tan x$: oppervlakte $\frac{\tan x}{2}$.

\textbf{3.} De driehoek ligt in de sector, en die in de grote driehoek:
$\frac{\sin x}{2} < \frac x2 < \frac{\tan x}{2}$, dus
$\sin x < x < \tan x$. Uit $\sin x < x$ volgt $\frac{\sin x}{x} < 1$; en uit
$x < \tan x = \frac{\sin x}{\cos x}$ volgt, na vermenigvuldigen met
$\frac{\cos x}{x} > 0$, dat $\cos x < \frac{\sin x}{x}$. Voor $x \to 0^+$ is
$\cos x \to 1$: ingesloten streeft $\frac{\sin x}{x}$ naar $1$; en
$\frac{\sin x}{x}$ is even, zodat de tweezijdige limiet $1$ is.

\textbf{4.} $\frac{1 - \cos x}{x^2} = \frac{1 - \cos^2 x}{x^2 (1 + \cos x)} =
\left(\frac{\sin x}{x}\right)^2 \frac{1}{1 + \cos x} \to 1 \times \frac12$.

\textbf{5.} $\sin'(0) = \lim_{h \to 0} \frac{\sin h - 0}{h} = 1$: het hele
afleiden van \omterm{def:g12:trigo:cossin}{sinus} en \omterm{def:g12:trigo:cossin}{cosinus} vloeit uit die ene limiet voort. In graden
gemeten zou de limiet $\frac{\pi}{180}$ zijn, en zou elke \omterm{def:g12:deriv:derivative}{afgeleide} die
constante meeslepen: \omterm{def:g11:trigo:radian}{radialen} zijn precies de eenheid die de analyse van de
trilling zuiver houdt.

\textbf{6.} $y' = -R\omega\sin(\omega t - \varphi)$ en
$y'' = -R\omega^2\cos(\omega t - \varphi) = -\omega^2 y$.

\textbf{7.} $\omega = 2$. $y(0) = A = 3$; $y'(0) = 2B = 8$, dus $B = 4$:
$y = 3\cos 2t + 4\sin 2t$.

\textbf{8.} $R = \sqrt{9 + 16} = 5$; periode $\frac{2\pi}{2} = \pi$; fase
$\varphi = \arctan\frac43 \approx 0.927$: $y = 5\cos(2t - 0.927)$.

\textbf{9.} $5\cos(2t - \varphi) = 0$ voor het eerst wanneer
$2t - \varphi = \frac\pi2$: $t = \frac{\pi/2 + 0.927}{2} \approx 1.25$ s.

\textbf{10.} Met $\sin\theta \approx \theta$ wordt
$\theta'' = -\frac gL \theta$: harmonisch met $\omega = \sqrt{\frac gL}$ en
periode $T = \frac{2\pi}{\omega} = 2\pi\sqrt{\frac Lg}$. Voor $L = 1$ is
$T \approx 2.006$ s. De periode hangt niet van de amplitude af (bij kleine
uitslagen) --- de isochronie van Galilei, het feit dat slingeruurwerken
mogelijk maakt.

\textbf{11.} $L = g\left(\frac{T}{2\pi}\right)^2 = \frac{9.81}{\pi^2}
\approx 0.994$ m: de secondeslinger komt ongeveer $6$ mm tekort op de
meridiaanmeter. De Academie koos de meridiaan (de zwaartekracht verschilt van
plaats tot plaats en verraadt de slinger); was $g$ een haartje groter
geweest, dan zou onze meter tikken.

\textbf{12.} $\cos\left(x + \frac\pi2\right) =
\cos x \cos\frac\pi2 - \sin x \sin\frac\pi2 = -\sin x$;
$\sin\left(x + \frac\pi2\right) = \cos x$. Dus geeft de \omterm{def:g12:trigo:cossin}{sinus} afleiden de
\omterm{def:g12:trigo:cossin}{sinus} een kwartdraai verder ($\cos$), en nog eens ($-\sin$), en nog eens, en
na vier stappen ben je thuis: afleiden draait de cirkel rond.

\textbf{13.} De twee uitwerkingen van $\cos(a \mp b)$ optellen geeft
$\cos(a-b) + \cos(a+b) = 2\cos a\cos b$. Met $a = b = x$:
$1 + \cos 2x = 2\cos^2 x$, dus $\cos^2 x = \frac{1 + \cos 2x}{2}$.

\textbf{14.} Met $p = \frac{p+q}{2} + \frac{p-q}{2}$ en $q$ als spiegelbeeld
geeft \omterm{def:g10:algebra:expand}{uitwerken} en optellen
$\sin p + \sin q = 2\sin\frac{p+q}{2}\cos\frac{p-q}{2}$. Voor $440$ en
$444$ Hz: een toon van $442$ Hz waarvan de luidheid gemoduleerd wordt door
$\cos(2\pi \cdot 2t)$ --- de luidheid piekt $4$ keer per seconde (de \omterm{def:g10:numbers:abs}{absolute
waarde} van de omhullende): vier zwevingen per seconde, die verdwijnen zodra
de stemvorken het eens zijn.

\textbf{15.} $\cos 3x = \cos 2x\cos x - \sin 2x \sin x =
(2\cos^2 x - 1)\cos x - 2\sin^2 x\cos x = 4\cos^3 x - 3\cos x$. In
$x = 20^\circ$ is $\cos 60^\circ = \frac12$, dus voldoet
$c = \cos 20^\circ$ aan $8c^3 - 6c - 1 = 0$ --- een derdegraadsvergelijking
zonder oplossing in vierkantswortels: $20^\circ$ construeren, dus $60^\circ$
in drieën delen, ligt buiten het bereik van passer en liniaal.

\textbf{16.} $R = \sqrt{9 + 27} = 6$; $\tan\varphi = \frac{3\sqrt3}{3} =
\sqrt3$: $\varphi = \frac\pi3$: $I(t) = 6\cos\left(100\pi t -
\frac\pi3\right)$: piek $6$ ampère, met een faseachterstand van $60^\circ$.

\textbf{17.} $\sin^2 t = \frac{1 - \cos 2t}{2}$: periode $\pi$. $\sin 2t$
heeft periode $\pi$ en $\sin 3t$ heeft $\frac{2\pi}{3}$: de som herhaalt zich
na het kleinste gemeenschappelijke veelvoud, $2\pi$.

\textbf{18.} Een trilling met periode $1$ binnen de krimpende omhullende
$\pm\eu^{-t/10}$: elke slag ongeveer $10\,\%$ lager. In $t = 10$ is de
omhullende $\eu^{-1} \approx 0.368$ --- opnieuw de constante van
\cref{pb:g12:exp:1}. \omterm{def:g10:algebra:equation}{Vergelijkingen} met zulke gedempte oplossingen
($y'' + a y' + b y = 0$) worden opgelost in het hoofdstuk over
differentiaalvergelijkingen.

\textbf{19.} Een kracht die evenredig met de uitwijking terugtrekt geeft
$y'' = -\omega^2 y$; en \cref{exo:g12:trigo:9} toont dat de oplossingen van
die \omterm{def:g10:algebra:equation}{vergelijking} \emph{precies} de combinaties van $\cos$ en $\sin$ zijn ---
bestaan door aanwijzing, uniciteit door de oefening: een trilling kan niet
anders dan sinusvormig zijn.

\textbf{20.} Eén limiet ($\frac{\sin x}{x} \to 1$, met ingesloten
oppervlakten), twee \omterm{def:g12:deriv:derivative}{afgeleiden} ($\sin' = \cos$, $\cos' = -\sin$), één
\omterm{def:g10:algebra:equation}{vergelijking} ($y'' = -\omega^2 y$), en een wereld van klokken, snaren,
stromen en getijden. Het sterretje: voor grote uitslagen is
$T = 4\sqrt{\frac Lg} \cdot \frac{\pi/2}{M(1, \cos(\theta_0/2))}$ --- het
\omterm{def:g12:seq:arith-geom}{rekenkundig-meetkundig} \omterm{def:g11:stat:mean}{gemiddelde} van \cref{pb:g12:seq:1} berekent de
elliptische integraal in een handvol iteraties: de dagboeknotitie van Gauss
en de slinger van Galilei, één formule van elkaar verwijderd.
\end{solution}
